\begin{lstlisting}[language=Python, numbers=left, basicstyle=\ttfamily]
clii_anchors = {
    "Cognitive": [
        "knowledge about climate change, e.g. Knowing that the melting of permafrost affects the Earth's climate or understanding the role of roof gardens in reducing carbon dioxide levels, notes UNCC:Learn",
        "understanding of environmental issues e.g. Recognizing that burning fossil fuels for power, cutting down forests, and manufacturing goods are major causes of climate change and that these actions lead to effects like rising temperatures and extreme weather events, according to the United Nations and the National Oceanic and Atmospheric Administration (NOAA).",
        "climate literacy and awareness, e.g. analyzing news articles or social media posts about climate change, discerning scientific consensus from misinformation, and understanding the uncertainties in scientific findings"
    ],
    "Affective": [
        "emotional response to climate change, e.g. This category includes a range of emotions such as eco-anxiety, a chronic fear of environmental doom; guilt about personal lifestyle choices; sadness and grief over the loss of ecosystems; and anger at inaction from governments and corporations. These negative emotions can be contrasted with hope and optimism, which can also motivate positive action.",
        "motivation to take action, e.g. The motivation to act on climate change is driven by a sense of duty and responsibility to future generations. It also relies on a person's efficacy beliefs, or the conviction that their actions can make a difference. Health and safety concerns, when climate change is framed as a direct threat to personal well-being, are also powerful motivators.",
        "concern and care for the environment, e.g. A person's level of concern is strongly linked to their emotional connection to nature and a deep sense of belonging to the natural world. Conversely, psychological distancing---the feeling that climate change is a problem for someone else---can reduce a person's concern. Finally, social norms and influence play a role, as people are often motivated to adopt pro-environmental behaviors to align with their social group. "
    ],
    "Behavioral": [
        "Actions to mitigate climate change: These encompass both individual and large-scale actions aimed at reducing or preventing greenhouse gas emissions. Examples include reducing personal energy consumption, adopting renewable energy sources, and switching to sustainable transport like walking, biking, or electric vehicles. It also involves lifestyle changes such as eating more plant-based foods, reducing food waste, and consciously minimizing consumption of goods like fast fashion and plastics.",
        "Community participation in sustainability: This involves the active and genuine involvement of people in addressing local problems and implementing projects to promote sustainable development. Participation can be driven by a motivation to improve one's community, building a sense of collective purpose. Strategies include community needs assessments, mobilization, and working with local authorities. It is crucial that participation is inclusive, involving women and other vulnerable groups, and that agencies or governments respect local knowledge and needs.",
        "Adoption of eco-friendly practices: The adoption of eco-friendly practices by individuals and organizations includes a range of actions that minimize environmental harm, reduce waste, and conserve energy. Examples involve using reusable items, recycling, upcycling, composting, and conserving water. In the corporate sector, this includes implementing strategies to reduce a company's carbon footprint, investing in sustainable supply chains, and offering eco-friendly products."
        ],
    "Institutional": [
        "policy development for climate adaptation: Effective policy-making for climate change adaptation involves national and sub-national governments developing strategies and measures to reduce vulnerability to climate impacts. This includes implementing large-scale measures like enhancing early warning systems, strengthening infrastructure, and developing strategies for managing water, land, and food systems. It also requires a 'whole-of-government' approach and integrating climate considerations into development plans.",
        "Governance and institutional support: This involves building robust institutional arrangements and governance models to address social problems like climate change and ensure sustainable development. Good governance is characterized by participatory behaviour, transparency, accountability, and coordination across different levels of government. Institutional support, often technical and strategic, is aimed at strengthening government departments to improve financial management, service delivery, and overall effectiveness.",
        "Capacity building in organizations: This is the process of helping individuals, organizations, and communities strengthen their skills, knowledge, and systems to achieve their goals. For organizations, this involves upskilling and reskilling employees, developing leadership, forming partnerships, and upgrading technology and processes. Ultimately, capacity building aims to transform behaviours and improve an organization's ability to achieve its mission."
    ]
}
\end{lstlisting}